%! Logarithmic Function Manipulation — Unit 5, Lesson 5.4 — Group Activity (Answer Key)
\documentclass[10pt]{article}
\usepackage{precalculus-article}
\usepackage{precalculus-key}   % pulls in -boxes; do NOT also load -boxes
\newcommand{\TallMath}[1]{$\displaystyle #1\rule[-1.4em]{0pt}{3.2em}$}

\begin{document}

\pageheader{Unit 5, Lesson 5.4}{Group Activity --- Answer Key}
\namepartnerperiod

\vspace{0.1in}

\begin{scenariobox}[Signal Power Analysis]{sky}
A radio engineer uses the formula $\text{dB} = 10 \cdot \log_{10}(P/P_0)$ to measure signal
power levels. Your group will simplify and rewrite logarithmic expressions to help her
analyze signals with different power levels.
\end{scenariobox}

\vspace{0.1in}

\begin{tcolorbox}[colback=white, colframe=black!40,
    title=\textbf{Tier R --- Remediate},
    fonttitle=\bfseries, arc=2mm, left=3mm, right=3mm, top=2mm, bottom=2mm]

\textbf{Apply the product and power properties to numerical expressions.}

\begin{enumerate}[leftmargin=1.8em, itemsep=10pt]

\item Rewrite each expression as a sum using the product property.

\vspace{4pt}
\begin{enumerate}[leftmargin=1.5em, itemsep=8pt, label=\alph*.]
  \item $\log_2(4 \cdot 8) = \log_2($ \ans{4} $) + \log_2($ \ans{8} $) = $ \ans{2} $+$ \ans{3} $=$ \ans{5}
  \item $\log_3(9 \cdot 27) = \log_3($ \ans{9} $) + \log_3($ \ans{27} $) = $ \ans{2} $+$ \ans{3} $=$ \ans{5}
  \item $\log_5(25 \cdot 125) = \log_5($ \ans{25} $) + \log_5($ \ans{125} $) = $ \ans{2} $+$ \ans{3} $=$ \ans{5}
\end{enumerate}

\vspace{4pt}
\item Apply the power property to rewrite each expression.

\vspace{4pt}
\begin{enumerate}[leftmargin=1.5em, itemsep=8pt, label=\alph*.]
  \item $\log_2(4^2) = $ \ans{2} $\cdot \log_2($ \ans{4} $) = $ \ans{2} $\cdot$ \ans{2} $=$ \ans{4}
  \item $\log_3(27^3) = $ \ans{3} $\cdot \log_3($ \ans{27} $) = $ \ans{3} $\cdot$ \ans{3} $=$ \ans{9}
\end{enumerate}

\end{enumerate}
\end{tcolorbox}

\vspace{0.1in}

\begin{tcolorbox}[colback=white, colframe=black!40,
    title=\textbf{Tier A --- Approaching Proficiency},
    fonttitle=\bfseries, arc=2mm, left=3mm, right=3mm, top=2mm, bottom=2mm]

\textbf{Rewrite algebraic expressions and apply change of base.}

\begin{enumerate}[leftmargin=1.8em, itemsep=10pt]

\item Rewrite $\log_2(32x)$ in the form $a + \log_2(x)$. What is the value of $a$?

\vspace{4pt}
$\log_2(32x) = \log_2($ \ans{32} $) + \log_2(x) = $ \ans{5} $+ \log_2(x)$

\vspace{2pt}
$a = $ \ans{5}. This represents a vertical \ans{translation up 5 units} of the graph of $y = \log_2(x)$.

\begin{teachernote}
EK 2.12.A.1: the horizontal dilation by factor 32 ($32x$ instead of $x$) is equivalent to
shifting the graph of $\log_2 x$ up by $\log_2 32 = 5$.
\end{teachernote}

\item Rewrite $\log_3(x^4)$ as a vertical dilation of $\log_3(x)$.

\vspace{4pt}
$\log_3(x^4) = $ \ans{$4\log_3(x)$}. The dilation factor is \ans{4}.

\item Use change of base to evaluate $\log_5(17)$ using common logarithms.

\vspace{4pt}
$\log_5(17) = \dfrac{\log({\color{keyred}\mathbf{17}})}{\log({\color{keyred}\mathbf{5}})} \approx $ \ans{1.760}

\item Show that $\log_2(x) = \dfrac{\ln(x)}{\ln(2)}$ using the change-of-base formula.

\ansline{By change of base with $a = e$: $\log_2(x) = \dfrac{\log_e(x)}{\log_e(2)} = \dfrac{\ln(x)}{\ln(2)}$.
This is a special case of EK 2.12.A.3 using the natural log as the intermediate base.}

\end{enumerate}
\end{tcolorbox}

\vspace{0.1in}

\begin{tcolorbox}[colback=white, colframe=black!40,
    title=\textbf{Tier E --- Extension},
    fonttitle=\bfseries, arc=2mm, left=3mm, right=3mm, top=2mm, bottom=2mm]

\textbf{Prove and justify log property connections to graph transformations.}

\begin{enumerate}[leftmargin=1.8em, itemsep=10pt]

\item Prove that $f(x) = \log_3(9x) = g(x) + 2$ where $g(x) = \log_3(x)$.

\ansline{By the product property: $\log_3(9x) = \log_3(9) + \log_3(x) = 2 + \log_3(x) = g(x) + 2$.
This is a vertical translation up 2 units; the horizontal dilation by 9 is equivalent to
shifting the output up by $\log_3(9) = 2$.}

\item Show that every $\log_b(x)$ is a vertical dilation of $\log_{10}(x)$.

\ansline{By change of base: $\log_b(x) = \dfrac{\log_{10}(x)}{\log_{10}(b)} = \dfrac{1}{\log_{10}(b)} \cdot \log_{10}(x)$.
This is $\log_{10}(x)$ multiplied by the constant $\dfrac{1}{\log_{10}(b)}$, which is a
vertical dilation. Since the dilation factor depends only on $b$ (not on $x$), every
logarithmic function is a vertical dilation of $\log_{10}(x)$.}

\item Simplify $\log_2(8x^3)$ using both properties with full justification.

\ansline{$\log_2(8x^3)$\\[3pt]
$= \log_2(8) + \log_2(x^3)$ \quad (product property)\\[3pt]
$= 3 + 3\log_2(x)$ \quad (since $\log_2(8) = 3$ and power property: $\log_2(x^3) = 3\log_2(x)$)\\[3pt]
$= 3 + 3\log_2(x)$}

\end{enumerate}
\end{tcolorbox}

\end{document}
